\documentclass[11pt]{article}
\usepackage[margin=1in]{geometry}
\usepackage[T1]{fontenc}
\usepackage[utf8]{inputenc}
\usepackage{lmodern}
\usepackage{microtype}
\usepackage{xcolor}
\usepackage{hyperref}
\usepackage{amsmath,amssymb}
\usepackage{booktabs,longtable,array}
\usepackage{enumitem}
\usepackage{listings}
\usepackage{fancyhdr}
\usepackage{titlesec}
\usepackage{graphicx}
\usepackage{float}
\usepackage{caption}
\usepackage{setspace}

\hypersetup{
  colorlinks=true,
  linkcolor=blue!50!black,
  urlcolor=blue!50!black,
  citecolor=blue!50!black,
  pdftitle={Integrated SAL Benchmark Code Guide},
  pdfauthor={OpenAI}
}

\pagestyle{fancy}
\fancyhf{}
\lhead{Integrated SAL Benchmark Code Guide}
\rhead{\thepage}
\renewcommand{\headrulewidth}{0.4pt}

\titleformat{\section}{\Large\bfseries\color{blue!45!black}}{\thesection}{0.6em}{}
\titleformat{\subsection}{\large\bfseries\color{blue!35!black}}{\thesubsection}{0.6em}{}

\lstdefinestyle{code}{
  basicstyle=\ttfamily\small,
  keywordstyle=\color{blue!60!black}\bfseries,
  commentstyle=\color{green!35!black},
  stringstyle=\color{red!50!black},
  showstringspaces=false,
  frame=single,
  rulecolor=\color{black!20},
  breaklines=true,
  columns=fullflexible,
  keepspaces=true,
  backgroundcolor=\color{black!02}
}

\setlist[itemize]{topsep=4pt,itemsep=2pt}
\setlist[enumerate]{topsep=4pt,itemsep=2pt}
\setlength{\parskip}{0.45em}
\setlength{\parindent}{0pt}

\begin{document}

\begin{titlepage}
\centering
\vspace*{2.2cm}

{\Huge\bfseries Integrated Sampler-Adaptive Learning Benchmark\\[0.4em]}
{\Large Code Guide and Design Notes\\[1.2em]}

{\large for \texttt{benchmark\_sal\_integrated\_ml\_physics\_qpu\_vs\_classical.py}}\\[2.0em]

\begin{minipage}{0.9\textwidth}
\centering
\setstretch{1.08}
This document explains, in detail, how the new benchmark code works. The code
combines sampler-adaptive learning ideas from the 2025 SAL paper with the
previous machine-learning benchmark and the earlier generalized transverse-field
Ising-model code. The result is one integrated program that can train RBMs,
masked RBMs, and semi-restricted Boltzmann machines using exact, Gibbs,
simulated-annealing, Langevin-SB, Pegasus-QPU, or Zephyr-QPU negative-phase
samplers, while also supporting physics-derived datasets from long-range 1D
TFIM, frustrated J1-J2 TFIM, and 2D Edwards-Anderson spin glass models.
\end{minipage}

\vfill

\begin{tabular}{>{\raggedright\arraybackslash}p{0.25\textwidth} p{0.62\textwidth}}
\toprule
Main benchmark script & \texttt{benchmark\_sal\_integrated\_ml\_physics\_qpu\_vs\_classical.py} \\
QPU helper & \texttt{qpu\_runtime\_safe\_sampler\_extended.py} \\
Primary training framework & SAL-style positive phase + sampler-specific negative phase \\
Extra physics metric & Variational energy of $\psi_\theta(s)=\sqrt{Q_\theta(s)}$ \\
\bottomrule
\end{tabular}

\vspace*{1.2cm}
\end{titlepage}

\section{What this program is trying to do}

The benchmark has two goals.

First, it compares several sampling backends inside the \emph{same} Boltzmann-machine
training loop:
\begin{itemize}
  \item exact marginalization,
  \item Gibbs sampling,
  \item simulated annealing,
  \item Langevin simulated bifurcation (LSB),
  \item a Pegasus-targeted D-Wave QPU backend,
  \item a Zephyr-targeted D-Wave QPU backend.
\end{itemize}

Second, it evaluates those samplers on both standard machine-learning datasets
and physics-derived datasets. For the physics-derived case, the code does not
stop at generative metrics such as KL divergence. It also interprets the trained
Boltzmann machine as a positive wavefunction and reports the corresponding
variational energy for the chosen Hamiltonian. This makes the script useful for
benchmarking both ``how well does the model learn the data distribution?'' and
``how physically good is the learned positive-amplitude state?''

\subsection{Why the SAL paper matters here}

The SAL paper contributes two ideas that are central to this code:
\begin{enumerate}
  \item a fast non-MCMC sampler, Langevin simulated bifurcation (LSB), and
  \item conditional expectation matching (CEM), which estimates an effective
  inverse temperature $\beta_{\mathrm{eff}}$ for samplers whose output temperature
  is not known a priori.
\end{enumerate}

The new code generalizes those ideas beyond the exact RBM/SRBM setup of the
paper in three directions:
\begin{itemize}
  \item it adds the \texttt{masked\_rbm} architecture,
  \item it allows non-LSB negative-phase samplers such as simulated annealing and
  real D-Wave QPUs to reuse the same CEM temperature-estimation machinery,
  \item it connects data-driven BM training to physics targets built from
  transverse-field Ising Hamiltonians.
\end{itemize}

\section{Files and their responsibilities}

The implementation is split into two files.

\subsection{Main script}

\texttt{benchmark\_sal\_integrated\_ml\_physics\_qpu\_vs\_classical.py} is the
orchestrator. It owns:
\begin{itemize}
  \item dataset generation,
  \item physical Hamiltonian generation,
  \item model construction,
  \item all training backends except the low-level Ocean glue,
  \item CEM fitting,
  \item the training loop,
  \item plotting,
  \item multiprocessing and output management.
\end{itemize}

\subsection{QPU helper}

\texttt{qpu\_runtime\_safe\_sampler\_extended.py} isolates Ocean- and
embedding-specific logic. It owns:
\begin{itemize}
  \item solver selection by topology or explicit solver name,
  \item embedding search and caching,
  \item support for disconnected and bias-only logical graphs,
  \item safe read-budget calculation with
  \texttt{estimate\_qpu\_access\_time()},
  \item repeated batched submissions under the runtime limit,
  \item return of expanded samples in a deterministic variable order.
\end{itemize}

The ``disconnected-graph'' improvement is especially important for SAL/CEM,
because the conditional hidden-only problems used by CEM often contain no
couplers at all; they are just collections of independent hidden spins with
shifted biases.

\section{Mathematical model used by the code}

\subsection{Energy and score conventions}

The code parameterizes a Boltzmann machine with visible spins
$v\in\{-1,+1\}^{N_v}$ and hidden spins $h\in\{-1,+1\}^{N_h}$ through
visible biases $a$, hidden biases $b$, hidden-visible couplings $W$, and
optional visible-visible couplings $L$.

Its unnormalized joint \emph{score} is
\begin{equation}
S(v,h)=a\cdot v+b\cdot h+h^T W v + \frac{1}{2} v^T L v.
\end{equation}

This is the sign-convention twin of the paper's energy
$E(v,h)= -S(v,h)$. Training therefore uses a gradient-ascent style
\emph{positive minus negative} update, whereas the paper writes the same physics
as gradient descent on the energy-form KL objective.

\subsection{Visible marginal at finite effective inverse temperature}

Given an effective inverse temperature $\beta$, the visible marginal is
\begin{equation}
Q_\beta(v) \propto \sum_h e^{\beta S(v,h)}
= \exp\left[\beta\left(a\cdot v + \frac12 v^T L v\right)
+ \sum_{j=1}^{N_h} \log 2\cosh\!\big(\beta[b_j + (Wv)_j]\big)\right].
\end{equation}

This expression is implemented by the function
\texttt{visible\_log\_unnorm\_beta()}. A key design choice is that the code uses
this exact finite-$\beta$ visible marginal everywhere it matters: in the exact
backend, in exact KL evaluation, in pseudo-log-likelihood evaluation, and in the
physics-side variational-energy calculation.

\subsection{Model families}

The same parameterization supports three architectures.

\begin{center}
\begin{tabular}{>{\raggedright\arraybackslash}p{0.2\textwidth} p{0.69\textwidth}}
\toprule
RBM & Dense hidden-visible bipartite couplings, no visible-visible edges. \\
Masked RBM & Same as RBM, but the hidden-visible matrix $W$ is sparsified by a random mask. \\
SRBM & Dense hidden-visible couplings plus a visible-visible mask $L$. The visible mask can be ring, dense, grid, or the interaction graph of the selected physics problem. \\
\bottomrule
\end{tabular}
\end{center}

\section{Datasets and physical targets}

\subsection{Machine-learning style datasets}

The code implements four purely data-distribution benchmarks.

\begin{itemize}
  \item \textbf{Bars and stripes}: all valid binary bar or stripe patterns on a
  small grid.
  \item \textbf{Parity}: all visible states with fixed parity.
  \item \textbf{Planted Ising}: a classical pairwise Ising distribution on a ring.
  \item \textbf{Three-spin}: an exact Boltzmann distribution of a random
  three-body spin-glass Hamiltonian, included because it mirrors the SAL paper's
  first synthetic rugged-landscape benchmark.
\end{itemize}

Each of these generators returns
\begin{itemize}
  \item the sampled training data,
  \item the exact support states when feasible,
  \item the exact target probabilities on that support.
\end{itemize}

That exact support is what makes it possible to evaluate KL divergence exactly
for small systems.

\subsection{Physics-derived dataset: \texttt{physics\_ground}}

For the physics mode, the code first builds a stoquastic transverse-field Ising
Hamiltonian and then computes the exact ground-state distribution in the
computational basis,
\begin{equation}
P_D(s)=|\psi_0(s)|^2.
\end{equation}

The supported Hamiltonians are:
\begin{align}
H_{\mathrm{long}} &= -\sum_{i<j} \frac{1}{r_{ij}^{\alpha}}\sigma_i^z\sigma_j^z
- \Gamma\sum_i \sigma_i^x,\\
H_{J_1J_2} &= -J_1\sum_i \sigma_i^z\sigma_{i+1}^z
- J_2\sum_i \sigma_i^z\sigma_{i+2}^z
- \Gamma\sum_i \sigma_i^x,\\
H_{\mathrm{EA}} &= -\sum_{\langle ij\rangle} J_{ij}\sigma_i^z\sigma_j^z
- \Gamma\sum_i \sigma_i^x,
\end{align}
with $J_{ij}=\pm 1$ in the Edwards-Anderson case.

The exact ground-state probabilities are obtained by sparse diagonalization of
the Hamiltonian matrix, using \texttt{eigsh()} on the stoquastic sparse operator.

\subsection{Why this physics mode is useful}

This design lets the code ask two distinct questions.

\begin{enumerate}
  \item Can a trained BM reproduce the exact ground-state probability
  distribution in the $\sigma^z$ basis?
  \item If that BM is converted into a positive wavefunction,
  is it also good in variational-energy terms?
\end{enumerate}

This is the bridge between ``energy-based generative modeling'' and ``neural
ansatz for many-body physics.''

\section{How the backends are implemented}

\subsection{Exact backend}

The exact backend enumerates all visible states and evaluates the finite-$\beta$
visible marginal exactly. It then computes
\begin{itemize}
  \item $\langle v_i\rangle$,
  \item $\langle h_j\rangle$,
  \item $\langle h_j v_i\rangle$,
  \item $\langle v_i v_j\rangle$,
\end{itemize}
by marginalization.

This backend is the ``gold standard'' negative phase whenever $N_v$ is small
enough that the support size $2^{N_v}$ remains manageable.

\subsection{Gibbs backend}

The Gibbs backend stores persistent visible and hidden chains. Each negative
phase performs:
\begin{enumerate}
  \item a blocked hidden update, because hidden units are conditionally
  independent given visible units,
  \item a sequential visible update, which is required when the model contains
  visible-visible couplings.
\end{enumerate}

For the RBM case, this reduces to the standard alternating blocked Gibbs sampler.
For the SRBM case, it becomes a proper semi-restricted sequential visible update.

\subsection{Simulated-annealing backend}

The simulated-annealing backend also stores persistent chains, but instead of
sampling from one fixed conditional distribution, it runs a geometric cooling
schedule from \texttt{sa\_temp\_start / beta\_submit} down to
\texttt{sa\_temp\_end / beta\_submit}. This makes its output distribution only an
\emph{effective} Boltzmann sampler, not an exact finite-temperature Markov chain.
That is exactly why CEM is useful: the positive phase should use the effective
$\beta_{\mathrm{eff}}$ inferred from the actual sampler output, not the nominal
cooling schedule.

\subsection{LSB backend}

The LSB backend is the direct algorithmic import from the SAL paper. It evolves
positions and momenta with
\begin{align}
Y &\leftarrow Y + \Delta\,(J\,\mathrm{sgn}(X)+f),\\
X &\leftarrow \mathrm{sgn}(X) + \Delta Y,
\end{align}
and after each step discretizes the positions and reinitializes the momenta from
a Gaussian distribution. In code this is implemented as:
\begin{enumerate}
  \item sample initial $S\in\{-1,+1\}^{N_v+N_h}$ and Gaussian $Y$,
  \item compute the force from the full BM graph,
  \item update $Y$ and then $X$,
  \item discretize back to spins,
  \item repeat for \texttt{lsb\_steps} iterations.
\end{enumerate}

The important conceptual point is that LSB does \emph{not} have a known target
inverse temperature. The code therefore always treats LSB as a backend requiring
CEM unless the user overrides that behavior with a fixed temperature.

\subsection{QPU backends}

The Pegasus and Zephyr backends use the same main ideas as the earlier QPU
benchmark code.

\begin{itemize}
  \item The BM is converted to an Ising model whose minimized energy is
  $-S(v,h)$.
  \item A fixed embedding is cached and reused.
  \item Complete bipartite RBMs use a biclique embedding fast path.
  \item SRBMs and other masked graphs fall back to generic minor-embedding.
  \item The helper estimates QPU access time before submission and splits the
  logical sample budget into safe child submissions.
\end{itemize}

The extended helper adds one more feature that the earlier code did not have:
if the conditional CEM problem has no couplers, the helper assigns isolated
logical variables to unused physical qubits one by one. This makes hidden-only
conditional sampling possible on the real QPU as well.

\section{Conditional expectation matching (CEM)}

\subsection{Why CEM is needed}

For exact Gibbs sampling at a chosen $\beta_{\mathrm{submit}}$, the effective
inverse temperature is known. For simulated annealing, LSB, and real QPU
sampling, the output distribution is only approximately Boltzmann and the actual
effective temperature can drift with model parameters and sampler settings.

The positive phase, however, needs hidden expectations of the form
\begin{equation}
\langle h_j\rangle_{v,\beta}=\tanh\!\Big(\beta\,[b_j+(Wv)_j]\Big).
\end{equation}

If the wrong $\beta$ is used here, the positive phase becomes inconsistent with
the negative phase. CEM fixes that mismatch.

\subsection{How CEM is implemented in the code}

The function \texttt{estimate\_beta\_eff()} performs the following steps.

\begin{enumerate}
  \item Choose one or more conditioning vectors $r$ from the training data, or
  from random visible states if the user requests that.
  \item Build the hidden-only conditional problem with fields
  \begin{equation}
  f_j(r)=b_j + (Wr)_j.
  \end{equation}
  \item Ask the current backend to generate conditional hidden samples under the
  same sampler settings used in the full negative phase.
  \item Compute empirical hidden means from those samples.
  \item Fit $\beta_{\mathrm{eff}}$ by minimizing
  \begin{equation}
  \frac{1}{KN_h}\sum_{k=1}^K\sum_{j=1}^{N_h}
  \left[m_{k,j}^{\mathrm{sample}}-\tanh\big(\beta f_j(r_k)\big)\right]^2.
  \end{equation}
\end{enumerate}

The minimization is one-dimensional and uses
\texttt{scipy.optimize.minimize\_scalar()} on a bounded interval.

\subsection{Estimator modes}

The user can select three modes.

\begin{itemize}
  \item \texttt{auto}: exact and Gibbs use the known submitted inverse
  temperature, all other samplers use CEM.
  \item \texttt{cem}: always run CEM.
  \item \texttt{fixed}: do not estimate anything; always use
  \texttt{fixed\_beta\_eff}.
\end{itemize}

The \texttt{auto} mode is generally the best default because it is faithful to
the SAL logic without paying unnecessary QPU cost when the temperature is known
analytically.

\section{Training loop, step by step}

The heart of the program is \texttt{run\_single\_experiment()}. Each epoch runs
in the following order.

\begin{enumerate}
  \item Estimate or reuse $\beta_{\mathrm{eff}}$.
  \item Select a minibatch if the user asked for minibatching; otherwise use the
  full dataset.
  \item Compute the positive-phase statistics using analytic hidden means at the
  current $\beta_{\mathrm{eff}}$.
  \item Generate negative-phase samples with the chosen backend.
  \item Update parameters by positive minus negative moments, with $L_2$
  regularization and clipping.
  \item Optionally evaluate pseudo-log-likelihood, reconstruction error, exact
  KL divergence, and physics variational energy.
  \item Save a detailed row in \texttt{history.csv}.
\end{enumerate}

In compact form, the update is
\begin{align}
a &\leftarrow a + \eta(\langle v\rangle_D - \langle v\rangle_{\mathrm{neg}})-\lambda a,\\
b &\leftarrow b + \eta(\langle h\rangle_D - \langle h\rangle_{\mathrm{neg}})-\lambda b,\\
W &\leftarrow W + \eta(\langle hv\rangle_D - \langle hv\rangle_{\mathrm{neg}})-\lambda W,\\
L &\leftarrow L + \eta(\langle vv\rangle_D - \langle vv\rangle_{\mathrm{neg}})-\lambda L,
\end{align}
where the data-side hidden expectations use $\beta_{\mathrm{eff}}$ and the
negative-phase moments come from the backend samples themselves.

\section{Why the physics-side variational energy is valid}

For a physics-derived dataset, the trained BM defines a visible probability
model $Q_\theta(s)$. The code then constructs the positive wavefunction
\begin{equation}
\psi_\theta(s)=\sqrt{Q_\theta(s)}.
\end{equation}

For the stoquastic transverse-field Ising models considered here, the local
energy for a basis state $s$ is
\begin{equation}
E_{\mathrm{loc}}(s)=E_z(s)-\Gamma\sum_i
\frac{\psi_\theta(s^{(i)})}{\psi_\theta(s)},
\end{equation}
where $E_z(s)$ is the diagonal Ising contribution and $s^{(i)}$ is the state with
spin $i$ flipped.

Because
\begin{equation}
\log \psi_\theta(s)=\frac12\log Q_\theta(s),
\end{equation}
these ratios can be evaluated exactly from the BM visible log probabilities.
That is what the function
\texttt{exact\_physics\_variational\_energy\_per\_spin()} does.

This is one of the strongest features of the code: a purely generative training
benchmark can be turned into a many-body physics benchmark without introducing a
second, unrelated ansatz implementation.

\section{QPU runtime safety and cost control}

\subsection{Why a dedicated helper is used}

Direct QPU use is expensive and failure-prone if runtime is not controlled. The
helper therefore performs the following procedure automatically.

\begin{enumerate}
  \item Determine or reuse a fixed embedding.
  \item Count the physical qubits required by that embedding.
  \item Query the selected solver for an estimated QPU access time.
  \item Binary-search the largest safe number of reads per child call under the
  chosen fraction of the solver runtime limit.
  \item Split the requested logical sample count across multiple child
  submissions.
\end{enumerate}

This is the same strategy that was already used in the earlier QPU benchmark,
but the new helper adds support for disconnected conditional CEM problems.

\subsection{Practical levers for reducing QPU time}

The code exposes several knobs that materially affect cost.

\begin{itemize}
  \item \texttt{negative\_samples}: the main cost driver.
  \item \texttt{cem\_samples}: the cost per CEM call.
  \item \texttt{cem\_period}: estimate $\beta_{\mathrm{eff}}$ every $k$ epochs
  instead of every epoch.
  \item \texttt{num\_spin\_reversal\_transforms}: affects batching overhead.
  \item \texttt{max\_runtime\_fraction}: determines how conservatively the code
  stays below the solver runtime limit.
\end{itemize}

For real QPU studies, it is usually wise to start with small values such as
\texttt{negative\_samples=128}, \texttt{cem\_samples=32}, and
\texttt{cem\_period=5} while debugging.

\section{Multiprocessing and output layout}

The top-level \texttt{main()} function builds one independent task per
(model, backend) pair. With \texttt{--processes > 1}, those tasks are dispatched
through a spawn-based multiprocessing pool. This keeps QPU and classical runs
isolated and avoids shared-state surprises.

The output directory structure is
\begin{lstlisting}[style=code]
output_dir/
  run_config.json
  dataset_info.json
  all_summaries.json
  figure_final_summary.png
  figure_<model>_training.png
  <dataset_name>/
    <model>/
      <backend>/
        history.csv
        summary.json
        final_params.npz
        best_params.npz
\end{lstlisting}

The most useful file for post-analysis is usually \texttt{history.csv}, because
it contains the epoch-by-epoch trace of
\begin{itemize}
  \item $\beta_{\mathrm{eff}}$,
  \item KL divergence,
  \item pseudo-log-likelihood,
  \item reconstruction error,
  \item variational energy if applicable,
  \item QPU timing and embedding diagnostics when relevant.
\end{itemize}

\section{Recommended usage patterns}

\subsection{Paper-style SAL sanity check}

Use the three-spin dataset with small visible dimension and compare exact,
Gibbs, SA, and LSB:

\begin{lstlisting}[style=code]
python benchmark_sal_integrated_ml_physics_qpu_vs_classical.py \
  --models rbm masked_rbm srbm \
  --samplers exact gibbs sa lsb \
  --dataset-source three_spin \
  --n-visible 10 \
  --n-hidden 5 \
  --epochs 100 \
  --negative-samples 512 \
  --output-dir runs_three_spin
\end{lstlisting}

\subsection{Physics benchmark with exact references}

\begin{lstlisting}[style=code]
python benchmark_sal_integrated_ml_physics_qpu_vs_classical.py \
  --models rbm srbm \
  --samplers exact gibbs sa \
  --dataset-source physics_ground \
  --physical-hamiltonian j1j2_1d \
  --j1j2-size 10 --j1 1.0 --j2 0.5 \
  --visible-graph physics \
  --transverse-field 1.0 \
  --epochs 80 \
  --negative-samples 512 \
  --output-dir runs_j1j2
\end{lstlisting}

Using \texttt{visible\_graph=physics} is a particularly sensible ablation here,
because it tests whether giving the SRBM the same graph as the Hamiltonian helps
relative to a generic ring or dense choice.

\subsection{Conservative first QPU run}

\begin{lstlisting}[style=code]
python benchmark_sal_integrated_ml_physics_qpu_vs_classical.py \
  --models rbm srbm \
  --samplers sa qpu_pegasus qpu_zephyr \
  --dataset-source physics_ground \
  --physical-hamiltonian long_range_1d \
  --long-range-size 8 \
  --transverse-field 0.8 \
  --epochs 30 \
  --negative-samples 128 \
  --cem-samples 32 \
  --cem-period 5 \
  --annealing-time 20 \
  --num-spin-reversal-transforms 2 \
  --output-dir runs_qpu_start
\end{lstlisting}

This is intentionally conservative. Once the pipeline behaves well, raise the
sample counts only for the final runs that go into plots.

\section{Extension points}

The code was written to be easy to extend.

\subsection{Add a new sampler}

To add another backend, implement a new subclass of \texttt{BackendBase} with
three methods:
\begin{itemize}
  \item \texttt{negative\_phase()},
  \item \texttt{conditional\_hidden\_means()},
  \item optionally \texttt{known\_beta\_eff()} if the sampler's temperature is
  analytically known.
\end{itemize}

This is enough to plug the new sampler into the same SAL training loop.

\subsection{Add a new physics Hamiltonian}

To add another stoquastic model, define a new builder that returns a
\texttt{PhysicalProblem} containing
\begin{itemize}
  \item the number of sites,
  \item the edge list,
  \item the coupling values,
  \item any longitudinal fields.
\end{itemize}

Then register it in \texttt{build\_physical\_problem()}.

\subsection{Add a new dataset}

To add a new classical target distribution, register it in \texttt{build\_dataset()}.
If it is small enough to enumerate exactly, provide support states and exact
probabilities. Doing so unlocks exact KL evaluation and exact model observable
computation automatically.

\section{Most important implementation choices}

If someone reads only one section of this guide, it should be this one.

\begin{enumerate}
  \item The code uses the exact finite-$\beta$ BM visible marginal, not a crude
  post-hoc division of the visible score by temperature.
  \item CEM is backend-agnostic: the same one-dimensional fit is reused for SA,
  LSB, and QPU backends.
  \item Physics targets are integrated at the dataset level, which means the
  same BM code can be scored by both generative and variational metrics.
  \item The QPU helper is robust to disconnected conditional problems, which is
  essential for practical SAL-style temperature estimation.
  \item Multiprocessing is done across full experiments rather than inside the
  sampler, which keeps the logic simple and avoids hidden shared state.
\end{enumerate}

\section{Short reference map of the main functions}

\begin{center}
\begin{longtable}{>{\ttfamily\raggedright\arraybackslash}p{0.34\textwidth} p{0.56\textwidth}}
\toprule
Function or class & Role \\
\midrule
\endfirsthead
\toprule
Function or class & Role \\
\midrule
\endhead
\texttt{build\_dataset} & Creates the selected dataset and its exact support metadata. \\
\texttt{build\_physical\_problem} & Constructs the chosen TFIM-style Hamiltonian graph. \\
\texttt{visible\_log\_unnorm\_beta} & Exact visible log unnormalized probability at inverse temperature $\beta$. \\
\texttt{positive\_phase\_stats} & Data-side statistics with analytic hidden means at $\beta_{\mathrm{eff}}$. \\
\texttt{ExactBackend} & Exact negative-phase marginalization. \\
\texttt{GibbsBackend} & Alternating blocked/sequential Gibbs sampler. \\
\texttt{SimulatedAnnealingBackend} & Cooling-schedule backend with unknown effective temperature. \\
\texttt{LangevinSBBackend} & LSB implementation inspired by the SAL paper. \\
\texttt{QPUBackend} & Pegasus or Zephyr backend using the runtime-safe helper. \\
\texttt{estimate\_beta\_eff} & Auto / CEM / fixed effective-temperature management. \\
\texttt{exact\_physics\_variational\_energy\_per\_spin} & Positive-wavefunction energy evaluator for physics datasets. \\
\texttt{run\_single\_experiment} & One complete training run for one model-backend pair. \\
\texttt{main} & Task generation, multiprocessing, figure generation, and overall orchestration. \\
\bottomrule
\end{longtable}
\end{center}

\section{Closing remarks}

This code is not just a direct port of the SAL paper, and it is not just a
minor revision of the earlier benchmark. It is a synthesis of both. The main
conceptual novelty is that one can now benchmark classical samplers, optimization
samplers, and real D-Wave QPUs inside one temperature-aware BM training loop,
and then score the result both as a generative model and, when a physical target
is used, as a positive variational state.

That is exactly the kind of structure you want when the real question is not
``which sampler wins on one number?'' but rather ``how do sampler quality,
effective temperature, embedding overhead, and physical-state fidelity interact
across model classes and hardware topologies?''

\end{document}
